\chapter{2018年全国II卷（文）}

\section{单选题}

\begin{problem}
\(\i(2+3\i)=\)（\quad）.

\choices
    {\(3-2\i\)}
    {\(3+2\i\)}
    {\(-3-2\i\)}
    {\(-3+2\i\)}

\end{problem}

\begin{problem}
已知集合 \(A=\{1,3,5,7\}\)，\(B=\{2,3,4,5\}\)，则 \(A\cap B=\)（\quad）.

\choices
    {\(\{3\}\)}
    {\(\{5\}\)}
    {\(\{3,5\}\)}
    {\(\{1,2,3,4,5,7\}\)}

\end{problem}

\begin{problem}
函数 \(f(x)=\dfrac{\e^x-\e^{-x}}{x^2}\) 的图象大致为

\choices
    {\choicebitmap[width=0.15\paperwidth]{img/2018/national_paper_2_liberal/q3_optA.png}}
    {\choicebitmap[width=0.15\paperwidth]{img/2018/national_paper_2_liberal/q3_optB.png}}
    {\choicebitmap[width=0.15\paperwidth]{img/2018/national_paper_2_liberal/q3_optC.png}}
    {\choicebitmap[width=0.15\paperwidth]{img/2018/national_paper_2_liberal/q3_optD.png}}

\end{problem}

\begin{problem}
已知向量 \(\bs a\)，\(\bs b\) 满足 \(\lvert \bs a\rvert=1\)，\(\bs a\cdot\bs b=-1\)，则 \(\bs a\cdot(2\bs a-\bs b)=\)（\quad）.

\choices
    {\(4\)}
    {\(3\)}
    {\(2\)}
    {\(0\)}

\end{problem}

\begin{problem}
从 2 名男同学和 3 名女同学中任选 2 人参加社区服务，则选中的 2 人都是女同学的概率为（\quad）.

\choices
    {\(0.6\)}
    {\(0.5\)}
    {\(0.4\)}
    {\(0.3\)}

\end{problem}

\begin{problem}
双曲线 \(\dfrac{x^2}{a^2}-\dfrac{y^2}{b^2}=1\)（\(a>0\)，\(b>0\)）的离心率为 \(\sqrt3\)，则其渐近线方程为（\quad）.

\choices
    {\(y=\pm\sqrt2\,x\)}
    {\(y=\pm\sqrt3\,x\)}
    {\(y=\pm\dfrac{\sqrt2}{2}x\)}
    {\(y=\pm\dfrac{\sqrt3}{2}x\)}

\end{problem}

\begin{problem}
在 \(\triangle ABC\) 中，\(\cos\dfrac{C}{2}=\dfrac{\sqrt5}{5}\)，\(BC=1\)，\(AC=5\)，则 \(AB=\)（\quad）.

\choices
    {\(4\sqrt2\)}
    {\(\sqrt{30}\)}
    {\(\sqrt{29}\)}
    {\(2\sqrt5\)}

\end{problem}

\begin{problem}
为计算 \(S=1-\dfrac12+\dfrac13-\dfrac14+\cdots+\dfrac1{99}-\dfrac1{100}\)，设计了如图的程序框图，则在空白框中应填入（\quad）.

\bitmapinlinefigure[0.65\linewidth][width=0.65\linewidth]{img/2018/national_paper_2_liberal/q8_fig1.png}

\choices
    {\(i=i+1\)}
    {\(i=i+2\)}
    {\(i=i+3\)}
    {\(i=i+4\)}

\end{problem}

\begin{problem}
在正方体 \(ABCD-A_1B_1C_1D_1\) 中，\(E\) 为棱 \(CC_1\) 的中点，则异面直线 \(AE\) 与 \(CD\) 所成角的正切值为（\quad）.

\choices
    {\(\dfrac{\sqrt2}{2}\)}
    {\(\dfrac{\sqrt3}{2}\)}
    {\(\dfrac{\sqrt5}{2}\)}
    {\(\dfrac{\sqrt7}{2}\)}

\end{problem}

\begin{problem}
若 \(f(x)=\cos x-\sin x\) 在 \([0,a]\) 是减函数，则 \(a\) 的最大值是（\quad）.

\choices
    {\(\dfrac{\pi}{4}\)}
    {\(\dfrac{\pi}{2}\)}
    {\(\dfrac{3\pi}{4}\)}
    {\(\pi\)}

\end{problem}

\begin{problem}
已知 \(F_1\)，\(F_2\) 是椭圆 \(C\) 的两个焦点，\(P\) 是 \(C\) 上的一点，若 \(PF_1\perp PF_2\)，且 \(\angle PF_2F_1=60^\circ\)，则 \(C\) 的离心率为（\quad）.

\choices
    {\(1-\dfrac{\sqrt3}{2}\)}
    {\(2-\sqrt3\)}
    {\(\dfrac{\sqrt3-1}{2}\)}
    {\(\sqrt3-1\)}

\end{problem}

\begin{problem}
已知 \(f(x)\) 是定义域为 \((-\infty,+\infty)\) 的奇函数，满足 \(f(1-x)=f(1+x)\). 若 \(f(1)=2\)，则 \(f(1)+f(2)+f(3)+\cdots+f(50)=\)（\quad）.

\choices
    {\(-50\)}
    {\(0\)}
    {\(2\)}
    {\(50\)}

\end{problem}

\section{填空题}

\begin{problem}
曲线 \(y=2\ln x\) 在点 \((1,0)\) 处的切线方程为\fillinblank{}.

\end{problem}

\begin{problem}
若变量 \(x\)，\(y\) 满足约束条件
\(\begin{cases}
x+2y-5\ge0,\\
x-2y+3\ge0,\\
x-5\le0
\end{cases}\)
则 \(z=x+y\) 的最大值为\fillinblank{}.

\end{problem}

\begin{problem}
已知 \(\tan\left(\alpha-\dfrac{5\pi}{4}\right)=\dfrac15\)，则 \(\tan\alpha=\fillinblank{}\).

\end{problem}

\begin{problem}
已知圆锥的顶点为 \(S\)，母线 \(SA\)，\(SB\) 互相垂直，\(SA\) 与圆锥底面所成角为 \(30^\circ\)，若 \(\triangle SAB\) 的面积为 \(8\)，则该圆锥的体积为\fillinblank{}.

\end{problem}

\section{解答题}


\begin{problem}
记 \(S_n\) 为等差数列 \(\{a_n\}\) 的前 \(n\) 项和，已知 \(a_1=-7\)，\(S_3=-15\).

\begin{enumerate}
    \item 求 \(\{a_n\}\) 的通项公式；
    \item 求 \(S_n\)，并求 \(S_n\) 的最小值.
\end{enumerate}

\end{problem}

\begin{problem}
如图是某地区 2000 年至 2016 年环境基础设施投资额 \(y\)（单位：亿元）的折线图.

为了预测该地区 2018 年的环境基础设施投资额，建立了 \(y\) 与时间变量 \(t\) 的两个线性回归模型. 根据 2000 年至 2016 年的数据（时间变量 \(t\) 的值依次为 \(1\)，\(2\)，\(\cdots\)，\(17\)）建立模型\(\circled{1}\)：\(\hat y=-30.4+13.5t\)；根据 2010 年至 2016 年的数据（时间变量 \(t\) 的值依次为 \(1\)，\(2\)，\(\cdots\)，\(7\)）建立模型\(\circled{2}\)：\(\hat y=99+17.5t\).

\begin{enumerate}
    \item 分别利用这两个模型，求该地区 2018 年的环境基础设施投资额的预测值；
    \item 你认为用哪个模型得到的预测值更可靠？并说明理由.
\end{enumerate}

\bitmapinlinefigure[0.88\linewidth][width=0.88\linewidth]{img/2018/national_paper_2_liberal/q18_fig1.png}

\end{problem}

\begin{problem}
如图，在三棱锥 \(P-ABC\) 中，\(AB=BC=2\sqrt2\)，\(PA=PB=PC=AC=4\)，\(O\) 为 \(AC\) 的中点.

\begin{enumerate}
    \item 证明：\(PO\perp\) 平面 \(ABC\)；
    \item 若点 \(M\) 在棱 \(BC\) 上，且 \(MC=2MB\)，求点 \(C\) 到平面 \(POM\) 的距离.
\end{enumerate}

\bitmapinlinefigure[0.42\linewidth][width=0.42\linewidth]{img/2018/national_paper_2_liberal/q19_fig1.png}

\end{problem}

\begin{problem}
设抛物线 \(C:y^2=4x\) 的焦点为 \(F\)，过 \(F\) 且斜率为 \(k\)（\(k>0\)）的直线 \(l\) 与 \(C\) 交于 \(A\)，\(B\) 两点，\(\lvert AB\rvert=8\).

\begin{enumerate}
    \item 求 \(l\) 的方程；
    \item 求过点 \(A\)，\(B\) 且与 \(C\) 的准线相切的圆的方程.
\end{enumerate}

\end{problem}

\begin{problem}
已知函数 \(f(x)=\dfrac13x^3-a(x^2+x+1)\).

\begin{enumerate}
    \item 若 \(a=3\)，求 \(f(x)\) 的单调区间；
    \item 证明：\(f(x)\) 只有一个零点.
\end{enumerate}

\end{problem}

\begin{problem}
在直角坐标系 \(xOy\) 中，曲线 \(C\) 的参数方程为
\(\begin{cases}
x=2\cos\theta,\\
y=4\sin\theta
\end{cases}\)
（\(\theta\) 为参数），直线 \(l\) 的参数方程为
\(\begin{cases}
x=1+t\cos\alpha,\\
y=2+t\sin\alpha
\end{cases}\)
（\(t\) 为参数）.

\begin{enumerate}
    \item 求 \(C\) 和 \(l\) 的直角坐标方程；
    \item 若曲线 \(C\) 截直线 \(l\) 所得线段的中点坐标为 \((1,2)\)，求 \(l\) 的斜率.
\end{enumerate}

\end{problem}

\begin{problem}
设函数 \(f(x)=5-\lvert x+a\rvert-\lvert x-2\rvert\).

\begin{enumerate}
    \item 当 \(a=1\) 时，求不等式 \(f(x)\ge0\) 的解集；
    \item 若 \(f(x)\le1\)，求 \(a\) 的取值范围.
\end{enumerate}

\end{problem}
